% ВЫВОДЫ СОГЛАСНО ПРАКТИКЕ!!!!!!!!!!!!!!


% В результате выполнения курсовой работы проведен анализ линейной цепи (ФНЧ третьего порядка) во временной и частотной областях.

% Вычислена передаточная функция $H(s)$. Найденные полюсы $$s_1 = -3.328$$ и $$s_{2,3} = -0.446 \pm j1.142$$ лежат в левой полуплоскости, что подтверждает устойчивость цепи. Наличие комплексно-сопряженных полюсов обуславливает колебательный характер переходного процесса.

% 2. Анализ частотных характеристик позволил определить полосу пропускания цепи, которая составила $$\Delta\omega \in [0; 1.06]$$.

% 3. Уравнения состояния решены аналитическим методом и численно с помощью алгоритма Эйлера. Практически полное совпадение графиков подтверждает корректность выбранного шага интегрирования $$\Delta t = 0.030$$.

% 4. Спектральный анализ показал, что степень искажения сигнала строго зависит от соотношения ширины его спектра и полосы пропускания цепи. Короткий импульс ($t_{\text{и2}}$), обладающий широким спектром ($5.99 \gg 1.06$), подвергся критическим искажениям формы. Длинный импульс ($t_{\text{и1}}$) исказился в меньшей степени.

% 5. Расчет реакции на периодическое воздействие с использованием отрезка ряда Фурье и метода в «замкнутой» форме показал высокую сходимость. Решение в замкнутой форме описывает поведение цепи абсолютно точно, тогда как спектральный метод демонстрирует осцилляции Гиббса в точках разрыва, но адекватно передает установившийся режим при достаточном количестве учитываемых гармоник.